\documentclass[11pt,a4paper]{article}

\usepackage[margin=2.5cm]{geometry}
\usepackage{amsmath,amssymb,amsthm}
\usepackage{algorithm}
\usepackage{algpseudocode}
\usepackage{booktabs}
\usepackage{graphicx}
\usepackage{hyperref}
\usepackage{xcolor}

\hypersetup{colorlinks=true,linkcolor=blue!50!black,citecolor=blue!50!black,
            urlcolor=blue!50!black}

\newtheorem{remark}{Remark}
\theoremstyle{definition}
\newtheorem*{definition*}{Definition}
\newcommand{\Bmat}{\mathbf{B}}
\newcommand{\Lmat}{\mathbf{L}}
\newcommand{\Dmat}{\mathbf{D}}
\newcommand{\Hmat}{\mathbf{H}}
\newcommand{\Rmat}{\mathbf{R}}
\newcommand{\Xmat}{\mathbf{X}}
\newcommand{\xvec}{\mathbf{x}}
\newcommand{\yvec}{\mathbf{y}}
\newcommand{\rvec}{\mathbf{r}}

\title{\textbf{Estimating a clustered localization radius\\
for the modified-Cholesky ensemble Kalman filter:\\
a combinatorial formulation}}

\author{Working note --- not a submission draft}
\date{\today}

\begin{document}
\maketitle

\begin{abstract}
\noindent
This note sets out, in one place, what the experiment suite computes and why it
is built the way it is. The localization radius of a modified-Cholesky ensemble
Kalman filter is treated as a vector of integers, one per cluster of grid
points, where the clusters come from ensemble-derived features rather than from
position. The resulting problem is combinatorial with $r_{\max}^{2K}$
configurations, and six metaheuristics are compared on it against two controls.
The quantity being minimized is computed from observations alone, by
withholding part of the observation lattice; an oracle that uses the true state
is run alongside as an upper bound on what any selection rule could achieve.
Several design decisions that look arbitrary are consequences of measurements
reported here, and two open problems are stated rather than hidden.
\end{abstract}

\section{What the problem is}

The analysis step of an ensemble Kalman filter combines a forecast ensemble
with observations. With an ensemble far smaller than the state dimension, the
sample covariance is rank deficient and carries spurious correlations between
distant variables. Localization removes them, and every localization scheme
depends on a radius of influence.

That radius is normally chosen by sweeping it against a known truth in a
synthetic experiment and then freezing it. The question here is whether it can
be \emph{estimated} instead --- from observations alone --- and whether it pays
to let it vary in space.

\subsection{The filter}

We use the modified-Cholesky estimator of the precision matrix
\cite{bickel2008,nino2018}. Let $\Xmat^b \in \mathbb{R}^{n \times N}$ be the
forecast ensemble, $\mathbf{\Delta X}$ its deviations from the mean. For each
component $i$ with predecessor set $P(i, r_i)$ --- the neighbours of $i$ within
radius $r_i$ whose index is lower --- solve the ridge regression
%
\begin{equation}
  \boldsymbol{\beta}_i
  = \big(\Xmat_{P}^{\top}\Xmat_{P} + \alpha \mathbf{I}\big)^{-1}
    \Xmat_{P}^{\top}\,\mathbf{\Delta X}_{i},
  \qquad
  \Xmat_{P} = \mathbf{\Delta X}_{P(i,r_i)}^{\top},
  \label{eq:ridge}
\end{equation}
%
and assemble
%
\begin{equation}
  \Lmat_{ii} = 1, \quad
  \Lmat_{i,P(i,r_i)} = -\boldsymbol{\beta}_i^{\top}, \quad
  \Dmat_{ii} = \operatorname{var}\!\big(\mathbf{\Delta X}_i - \Xmat_P \boldsymbol{\beta}_i\big)^{-1},
  \qquad
  \widehat{\Bmat}^{-1}(\rvec) = \Lmat^{\top}\Dmat\,\Lmat .
  \label{eq:precision}
\end{equation}
%
The analysis then solves
$\big(\widehat{\Bmat}^{-1} + \Hmat^{\top}\Rmat^{-1}\Hmat\big)\,
 \delta\xvec = \Hmat^{\top}\Rmat^{-1}(\yvec - \Hmat\xvec^b)$.

Two properties of \eqref{eq:precision} drive everything that follows.

\begin{remark}[Row separability]
Row $i$ of $\Lmat$ and entry $\Dmat_{ii}$ depend on $r_i$ \emph{and on nothing
else}. Changing one component of the radius vector therefore leaves every other
row untouched.
\end{remark}

\begin{remark}[Positive definiteness is structural]
$\widehat{\Bmat}^{-1}$ is positive definite by construction for any $\rvec$,
because it has the form $\Lmat^{\top}\Dmat\Lmat$ with $\Dmat$ diagonal and
positive. Truncating an already-formed precision matrix does not preserve this:
in a companion measurement, masking entries beyond a given distance produced an
indefinite matrix in $44$ of $44$ cases tested, with minimum eigenvalues as
negative as $-2.8\times 10^{4}$.
\end{remark}

\section{The decision variable}

\subsection{Clusters, not grid points}

Giving every component its own radius means $n$ parameters --- here $n = 18818$
--- fitted against a few hundred observations. That does not fail for want of a
better search. In a companion experiment on a smaller system, an oracle with
full access to the truth reduced the analysis error by $59\%$ on the cycle it
optimized, and the same radius field made the error $22\%$ \emph{worse} on the
following cycles, performing worse than a random shuffle of its own values. The
structure it found was the cycle's noise.

We therefore group the grid points of each field into $K$ clusters and estimate
one radius per cluster. The features are computed from the forecast ensemble
alone --- the log background variance and the ensemble correlation at lags one
and two, averaged over the four directions --- so the grouping is admissible in
the same sense as the criterion. Clustering is by $k$-means on those features.

\begin{remark}[The clusters need not be spatially contiguous]
Points are grouped by how their ensemble behaves, not by where they sit. Two
points at opposite corners of the domain can share a radius. One consequence
matters for the search: a single decision variable controls grid points
scattered over the whole domain, and the variables are close to independent of
one another.
\end{remark}

\subsection{An integer vector}

With two fields, the decision variable is
%
\begin{equation}
  \boldsymbol{\theta}
  = (\theta_1, \dots, \theta_K,\ \theta_{K+1}, \dots, \theta_{2K})
  \in \{1, \dots, r_{\max}\}^{2K},
\end{equation}
%
the first $K$ entries the radii of the clusters of $q$ and the next $K$ those of
$\psi$. The expansion to the $n$ per-component radii is the deterministic map
$r_j = \theta_{c(j)}$, with $c(j)$ the cluster of point $j$; it is not part of
the search. With $K = 4$ and $r_{\max} = 12$ the space has $12^{8} \approx
4.3 \times 10^{8}$ elements.

Nothing in the search produces a real-valued solution and nothing is rounded.

\section{The objective}

\subsection{Why observations must be withheld}

Scoring the analysis against the observations it assimilated is circular. In a
direct measurement, that objective fell monotonically as the radius shrank,
reaching its minimum at the smallest admissible value: a short radius makes each
point fit its own observation and the residual goes to zero. The criterion
rewards copying the data.

So the observation set $\mathcal{O}$ is split, $\mathcal{O} = \mathcal{T} \cup
\mathcal{V}$, the analysis is computed from $\mathcal{T}$ alone, and the result
is scored at $\mathcal{V}$:
%
\begin{equation}
  J(\boldsymbol{\theta})
  = \frac{1}{|\mathcal{V}|}
    \big\| \yvec_{\mathcal{V}}
           - \Hmat_{\mathcal{V}}\,\xvec^{a}(\boldsymbol{\theta}; \mathcal{T})
    \big\|^{2}_{\Rmat^{-1}_{\mathcal{V}}} .
  \label{eq:objective}
\end{equation}
%
No reference trajectory enters \eqref{eq:objective}.

\subsection{One partition or many}

A single partition already removes the circularity. Repeating it over $B$ draws
reduces variance, and the two natural aggregations differ:

\begin{itemize}
  \item \textbf{Averaging} $J$ over the $B$ partitions and minimizing the mean.
  \item \textbf{Voting}: minimize each realization separately and take, for each
        component, the mode of the $B$ winning values. Because the radii are
        integers the mode needs no binning.
\end{itemize}

Voting is the more robust of the two when the landscape is flat, since each
vote only has to fall in the right place rather than locate an exact minimum.
It is also the more expensive: $B$ optimizations instead of one.

\begin{remark}[A diagnostic that comes free]
The agreement among votes --- how many of the $B$ winners a component's mode
collects --- says whether the structure is identified or whether noise is being
averaged. In the companion experiment the agreement was $7$ out of $30$, which
is what a null result looks like.
\end{remark}

\subsection{The oracle}

Alongside \eqref{eq:objective} we run
$J_{\text{oracle}}(\boldsymbol{\theta}) = \|\xvec^{a}(\boldsymbol{\theta})
- \xvec^{\text{true}}\|$, which no operational system can compute. Its purpose
is to bound what any selection rule could achieve. If a clustered radius does
not beat a uniform one \emph{with the truth in hand}, there is nothing to
estimate and the negative result is clean.

\section{The searches}

All six share one neighbourhood, so that a comparison between them measures the
acceptance rule and not the move.

\begin{definition*}[Neighbourhood]
$\mathcal{N}(\boldsymbol{\theta})$ is the set of vectors that differ from
$\boldsymbol{\theta}$ in exactly one component. With a step $s$, the new value
is drawn from $\{\theta_j - s, \dots, \theta_j + s\} \cap [1, r_{\max}]$;
with $s = \infty$, from the whole range.
\end{definition*}

\begin{algorithm}[t]
\caption{Tabu search on the integer radius vector}
\begin{algorithmic}[1]
\State $\boldsymbol{\theta} \gets$ random vector in $\{1,\dots,r_{\max}\}^{2K}$;
       $\boldsymbol{\theta}^\star \gets \boldsymbol{\theta}$
\State $T \gets \emptyset$ \Comment{tabu list of (component, value) pairs}
\While{budget remains}
  \State $\mathcal{C} \gets \{\,\boldsymbol{\theta}' \in
         \mathcal{N}(\boldsymbol{\theta}) :
         (j, \theta'_j) \notin T \ \text{or}\
         J(\boldsymbol{\theta}') < J(\boldsymbol{\theta}^\star)\,\}$
         \Comment{aspiration}
  \State $\boldsymbol{\theta} \gets \arg\min_{\boldsymbol{\theta}'
         \in \mathcal{C}} J(\boldsymbol{\theta}')$
  \State add $(j, \theta_j)$ to $T$ for $\tau$ iterations
  \If{$J(\boldsymbol{\theta}) < J(\boldsymbol{\theta}^\star)$}
      $\boldsymbol{\theta}^\star \gets \boldsymbol{\theta}$
  \EndIf
\EndWhile
\end{algorithmic}
\end{algorithm}

\begin{algorithm}[t]
\caption{Discrete flower pollination}
\begin{algorithmic}[1]
\State initialize a population of $S$ integer vectors; $\boldsymbol{\theta}^\star
       \gets$ best of them
\While{budget remains}
  \For{each individual $\boldsymbol{\theta}^{(s)}$}
    \If{$\mathrm{rand} > p$} \Comment{global pollination}
      \State $m \gets$ a L\'evy draw mapped to $\{1,\dots,2K\}$
      \State copy $m$ randomly chosen components from
             $\boldsymbol{\theta}^\star$ into $\boldsymbol{\theta}^{(s)}$
    \Else \Comment{local pollination}
      \State pick $u, v$ at random; let $\mathcal{D} =
             \{j : \theta^{(u)}_j \neq \theta^{(v)}_j\}$
      \State copy a random subset of $\mathcal{D}$ from
             $\boldsymbol{\theta}^{(u)}$ into $\boldsymbol{\theta}^{(s)}$
    \EndIf
    \State accept if improved; update $\boldsymbol{\theta}^\star$
  \EndFor
\EndWhile
\end{algorithmic}
\end{algorithm}

\noindent
The two moves are the discrete reading of the continuous algorithm
\cite{yang2012}. The L\'evy step does not produce a position; it produces
\emph{how many} components a move copies, so its heavy tail gives mostly small
moves with occasional large ones, which is the property the flight contributes.
Following the reference implementation, the global step is taken when
$\mathrm{rand} > p$, so $p = 0.8$ means global pollination one time in five.

\medskip
\noindent
The remaining four:

\begin{description}
  \item[Simulated annealing] the same neighbourhood as tabu, accepting a
    worsening move with probability $\exp(-\Delta/T)$; the step shrinks with
    temperature, so it is a tabu search with an unrestricted but cooling
    neighbourhood rather than a different algorithm.
  \item[Genetic algorithm] uniform crossover on whole components --- integers
    have no meaningful midpoint --- with one-component mutation and tournament
    selection.
  \item[Ant colony] a pheromone table $\tau \in \mathbb{R}^{2K \times r_{\max}}$
    over (component, value) pairs. Each ant builds a whole vector by choosing a
    value per component independently with probability
    $\propto \tau_{j,v}^{\,a}$; the best ants reinforce, everything evaporates.
    \emph{This is the search the parameterization suits best}: because the
    clusters are not spatially contiguous their radii are close to independent,
    and the pheromone accumulates evidence about each one separately, so an ant
    can combine the good value of one cluster with that of another even though
    the two have never appeared together.
  \item[Firefly] attraction becomes a probability of copying: each component of
    a brighter individual is copied with probability $\beta$, decaying with the
    Hamming distance between the two vectors.
\end{description}

\noindent
Two \textbf{controls}, which are not competitors but the calibration of the
comparison: uniform random sampling at the same budget, and an exhaustive sweep
where the space is small enough to enumerate. When the sweep completes, every
other search is scored as a gap to the true optimum rather than against its
rivals.

\section{Cost, and what it forces}

\subsection{Row caching}

Row separability has a direct consequence. The pair $(i, r_i)$ determines row
$i$ of $\Lmat$ and entry $\Dmat_{ii}$ completely, so they can be cached. A move
that changes one cluster then recomputes only that cluster's rows, and a move
that revisits a radius recomputes nothing.

Profiling one analysis at $n = 4802$: $7.8$\,s in total, $7.2$ in the precision,
and $5.1$ of those inside \texttt{scikit-learn} --- $4800$ calls to
\texttt{Ridge.fit}, one per grid point, of which $2.1$\,s were input validation.
Each regression has four predecessors and forty members and takes microseconds
to solve directly. Solving \eqref{eq:ridge} in closed form and caching by row
takes a warm rebuild from $3.48$\,s to $0.126$\,s, a factor of $28$.

\begin{remark}[An argument for the trajectory methods]
A move that changes one component costs a fraction of a full evaluation, while
a move that changes the whole vector costs all of it. This is a property of the
problem, not of the implementation, and it favours tabu search and annealing
over the population methods in a way that budget-counted comparisons at equal
\emph{evaluations} do not show.
\end{remark}

\subsection{What still does not fit}

At $n = 18818$ with $N = 40$, one analysis costs $11$\,s. Optimizing inside
every assimilation cycle with $B$ partitions and a budget of $P$ evaluations
costs $B \cdot P \cdot 11$ seconds per cycle: with $B = 10$, $P = 40$ and $40$
cycles, about $49$ hours per run. The options are to optimize every $M$ cycles
rather than every cycle --- defensible, since the radius is supposed to vary
slowly --- to reduce $B$, or to update the analysis solve under a low-rank
change instead of re-solving. The last has not been implemented.

\section{The testbed}

The model is the $1.5$-layer quasi-geostrophic system of Sakov and
Oke~\cite{sakov2008},
%
\begin{equation}
  q_t = -\psi_x - \varepsilon J(\psi, q) - A\,\triangle^{3}\psi
        + 2\pi \sin(2\pi y), \qquad q = \triangle\psi - F\psi,
\end{equation}
%
on the unit square with $\psi = \triangle\psi = \triangle^2\psi = 0$ on the
boundary, $F = 1600$, $\varepsilon = 10^{-5}$. It was chosen because Lorenz-96
cannot answer the question: with a uniform grid, uniform forcing and identical
observations everywhere, no grid point has any reason to want a different radius
from its neighbour, so a negative result there says nothing about the idea. The
QG model has two fields with genuinely different correlation scales --- $q$
carries fine filaments while $\psi$ solves a Helmholtz problem on $q$ and is
smooth --- and a flow that is not homogeneous.

Three settings are consequences of measurements rather than choices.

\paragraph{The dissipation is ten times the model default.}
At $\mathrm{rkh2} = 10^{-12}$ the norm of $q$ grows without saturating, a factor
of fourteen between $t = 500$ and $t = 8000$, so there is no stationary regime
and no climatology to sample. At $10^{-11}$ it settles: $1.107 \times 10^{4}$ at
$t = 20000$ against $1.222 \times 10^{4}$ at $t = 60000$. Sakov and Oke report
the same, raising the dissipation by ten and reducing the step from $1.5$ to
$1.25$ for a stable assimilating system.

\paragraph{The ensemble is climatological.}
One long run, snapshots along it, members and truth drawn at random from that
set. Perturbing a state and propagating does not work here, and the reason was
measured: a perturbation is mostly energy off the attractor and the
hyperviscosity removes it. Starting from $30\%$ of climatology the background
error falls to $0.068$ after $120$ time units, a factor of five, and raising the
amplitude does not compensate because the propagation eats it. Only after about
$250$ units does the surviving component grow, reaching $0.872$ at $t = 1000$ ---
exactly what a climatological ensemble has from the start.

\paragraph{Observations sit on a regular lattice.}
With a random subset the spacing fluctuates, so some local domains hold several
observations and others none, and the analysis at a point with an empty domain
is just the background. Measured at $2\%$ random density, the analysis left $q$
unchanged to four decimal places at \emph{every} radius, which looks like a
broken filter and is not one.

\section{The experiments}

\begin{table}[h]
\centering
\begin{tabular}{@{}lll@{}}
\toprule
Identifier & Question & Output \\
\midrule
\textsc{exp-00-setup} & Is the testbed what it claims to be?
  & climatology, ensemble, lattice \\
\textsc{exp-01-regime} & Where is the radius a two-sided decision?
  & optimum vs.\ observation density \\
\textsc{exp-02-clusters} & Does a clustered radius pay, and which search finds it?
  & gain over uniform; search comparison \\
\bottomrule
\end{tabular}
\end{table}

\noindent
\textsc{exp-01} sweeps the observation lattice spacing against a uniform radius
and reports, per density, whether the optimum falls in the \emph{interior} of
the admissible range. An optimum at the boundary means the configuration has
not reached the regime where the radius trades one error against another, and
no estimation experiment run there would mean anything.

\textsc{exp-02} runs the six searches and the two controls at equal budget,
under both \eqref{eq:objective} and the oracle, for $K \in \{1,2,3,4,6\}$.

\section{Open problems}

\paragraph{The precision matrices do not yet agree to machine precision.}
On a single row the closed-form ridge reproduces \texttt{scikit-learn} to
$2.4 \times 10^{-16}$. After the boundary components were excluded from
estimation --- $q$ vanishes there, so $1/\mathrm{var}$ is infinite, and
\texttt{pyteda} produces $290$ infinities at $n = 4802$ --- the assembled
matrices differ by tens of percent, with the largest discrepancies on diagonal
entries of order $10^{13}$. Whether that is confined to those degenerate rows or
whether excluding the boundary changes the interior regressions has not been
traced. The equivalence test is marked expected-to-fail rather than relaxed
until it passes.

\paragraph{Simulated annealing underperforms random sampling.}
On a separable quadratic over four components, averaged across six seeds, tabu
reaches the exact optimum every time, random sampling reaches $3.0$, and
annealing reaches $4.0$; slower cooling makes it worse. It spends too much of a
small budget on temperature levels. It is left in the comparison as it
performs.

\paragraph{The state is normalized per field, and that is an implementation
decision.} Each field is divided by its climatological spread before the
analysis and restored after. The reason is the observation error: \texttt{pyteda}'s
LETKF reads a single scalar variance, so with $q$ of order $10^{3}$ and $\psi$
of order $1$ no isotropic error suits both. It is not part of the
modified-Cholesky method and can be removed by passing a heterogeneous $\Rmat$.

\begin{thebibliography}{9}
\bibitem{bickel2008} P.~J. Bickel and E.~Levina.
  Regularized estimation of large covariance matrices.
  \emph{The Annals of Statistics}, 36(1):199--227, 2008.
\bibitem{nino2018} E.~D. Nino-Ruiz, A.~Sandu, and X.~Deng.
  An ensemble Kalman filter implementation based on modified Cholesky
  decomposition for inverse covariance matrix estimation.
  \emph{SIAM Journal on Scientific Computing}, 40(2):A867--A886, 2018.
\bibitem{sakov2008} P.~Sakov and P.~R. Oke.
  A deterministic formulation of the ensemble Kalman filter: an alternative to
  ensemble square root filters. \emph{Tellus A}, 60(2):361--371, 2008.
\bibitem{yang2012} X.-S. Yang.
  Flower pollination algorithm for global optimization.
  In \emph{Unconventional Computation and Natural Computation}, LNCS 7445,
  pages 240--249, Springer, 2012.
\end{thebibliography}

\end{document}
